\chapter{Borrowed results}\label{ch:borrowed}
Results taken from the literature. Each names its paper, its library key in
\texttt{lib/literature/}, its theorem number when known, and its entry in
\texttt{sink/formalizing/facts.md}. The thinking spaces add to this chapter.

\begin{borrowed}[Tests and data processing]\blabel{fact-test-fact-data-processing}\status{borrowed}
Let $P_0,P_1$ be laws on one measurable space and $\gamma\ge1$. The best advantage
of a test, $\sup_E\,(P_1(E)-P_0(E))$ over measurable $E$, equals $\TV(P_0,P_1)$. For
every Markov kernel $K$, $\Hg(KP_0\|KP_1)\le\Hg(P_0\|P_1)$, where
$\Hg(P\|Q)=\sup_E\,(P(E)-\gamma Q(E))$. The divergence $\Hg$ is jointly convex
under mixtures.
\notation{$\Hg(P\|Q)$}{hockey-stick divergence, $\sup_E (P(E)-\gamma Q(E))$}
\notation{$\TV$}{total variation, $\Hg$ at $\gamma=1$}
\source{Standard: Neyman and Pearson for the test, data processing for kernels.}
\end{borrowed}

\begin{borrowed}[Generative leap]\blabel{fact-generative-leap-multi-index-rate}\status{borrowed}
For Gaussian multi-index models with generative leap exponent $k^\star$, detecting the
index space takes $n=\Theta(d^{\max(1,k^\star/2)})$ samples. The lower bound holds
for low-degree algorithms, and an efficient algorithm meets it.
\notation{$k^\star$}{generative leap exponent of a multi-index model}
\source{Damian, Lee and Bruna, NeurIPS 2025, library key \texttt{DLB25}, theorem
number to be filled when the paper is read. Facts:
\slug{fact-generative-leap-multi-index-rate}, and
\slug{fact-generative-exponent-measures-detection}, which records that the bound
measures detection, not estimation.}
\end{borrowed}

\section{From the space oracle-and-measurement}

\begin{borrowed}[LiRA prices membership inference in shadow runs]\blabel{fact-lira-priced-in-training-runs}\status{borrowed}
The likelihood-ratio attack trains $N$ shadow models on random halves of a data
pool, so each target lies in $N/2$ of them, with $N=256$ and $64$ on ImageNet, and
one bank serves every target. Its true-positive rate at low false-positive rate
saturates near $64$ shadow models.
\source{Carlini, Chien, Nasr, Song, Terzis and Tramèr, IEEE S\&P 2022, library key
\texttt{CCNSTT22}, \S IV-C (Algorithm 1), \S V-A with Fig.~5, \S VI-B with Fig.~9.
Empirical. Accepted by argon in round 00.}
\fullproof{fact-lira-priced-in-training-runs}{lab/theory/spaces/oracle-and-measurement/round-00/builder.md}
\end{borrowed}

\begin{borrowed}[Sample price of membership against a noisy mean]\blabel{fact-hsu-membership-sample-gap}\status{borrowed}
Against the noisy empirical mean over all Gaussians, no membership test with $m$
auxiliary samples succeeds with probability $0.51$ on each side when
$d>c_1(n+n^2\rho^2)$ and $m\le c_2(n+n^2\rho^2)$. A matching attack exists against
symmetric estimators at the critical dimension. With known identity covariance no
test exists for $d<c(n+n^2\rho^2+n^2/m)$.
\source{Haghifam, Smith and Ullman, arXiv 2508.19458 v1 (FOCS 2026 to appear),
library key \texttt{HSU25}, Theorems 3.1, 3.3 and 4.1, Definitions 2.1, 2.4 and 3.2.
Accepted by argon in round 02.}
\fullproof{fact-hsu-membership-sample-gap}{lab/theory/spaces/oracle-and-measurement/round-02/builder.md}
\end{borrowed}

\begin{borrowed}[A small seed space breaks the privacy claim]\blabel{fact-small-seed-space-breaks-dp}\status{borrowed}
A DP-SGD implementation whose noise seed takes 100 values is detected as violating
its claimed $\varepsilon=1.27$ by the Gaussian-DP audit and missed by
$(\varepsilon,\delta)$ audits.
\source{Nasr, Hayes, Steinke, Balle, Tramèr, Jagielski, Carlini and Terzis, USENIX
Security 2023, library key \texttt{NHSBTJCT23}, \S 6.3, Violation 2, Fig.~8. Accepted
by argon in round 02.}
\fullproof{fact-small-seed-space-breaks-dp}{lab/theory/spaces/oracle-and-measurement/round-02/builder.md}
\end{borrowed}

\begin{borrowed}[Uncontrolled reruns at one seed diverge]\blabel{fact-rerun-reproduction-error}\status{borrowed}
Whole GPU training reruns at one seed end at an $\ell_2$ distance comparable to the
distance between independently trained models, while a single step nearly
reproduces.
\source{Jia, Yaghini, Choquette-Choo, Dullerud, Thudi, Chandrasekaran and Papernot,
IEEE S\&P 2021, library key \texttt{JYCDTCP21}, Theorem 1 and Table I. Empirical.
Accepted by astra in round 01.}
\fullproof{fact-rerun-reproduction-error}{lab/theory/spaces/oracle-and-measurement/round-01/builder.md}
\end{borrowed}

\begin{borrowed}[Controlled nondeterminism replicates reruns exactly]\blabel{fact-controlled-nondeterminism-exact-replication}\status{borrowed}
With deterministic settings reruns are identical on one GPU type and diverge across
types. Logging the rounding decisions replicates training exactly across three GPU
types at 1.2 to 1.7 times the training time.
\source{Srivastava, Arora and Boneh, NeurIPS 2024, library key \texttt{SAB24}, \S 6.
Empirical. Accepted by astra in round 01.}
\fullproof{fact-controlled-nondeterminism-exact-replication}{lab/theory/spaces/oracle-and-measurement/round-01/builder.md}
\end{borrowed}

\begin{borrowed}[A prover that picks the batches passes replay]\blabel{fact-unlearning-verification-fragile}\status{borrowed}
A proof of retraining, valid when every step replays within $\varepsilon$ and no batch
meets the unlearned set, is passed at $\varepsilon=0$ by retraining with each unlearned
point replaced by its nearest same-class neighbour, and at $\varepsilon>0$ by forging from
the logged proof in at most one run's time.
\source{Zhang, Chen, Shen and Li, ICML 2024, library
key \texttt{ZCSL24}, Definition 4.1, Propositions 4.2 to 4.5. Accepted by argon in round 02.}
\fullproof{fact-unlearning-verification-fragile}{lab/theory/spaces/oracle-and-measurement/round-02/builder.md}
\end{borrowed}

\section{From the space computational-limits}

\begin{borrowed}[SQ bounds for planted biclique]\blabel{fact-fgrvx17-planted-clique-sq-bound}\status{borrowed}
For every constant $\delta>0$, every $k\le n^{1/2-\delta}$ and every $r>0$, detecting
or finding a planted $k$-biclique with probability $2/3$ needs $n^{\Omega(\log r)}$
queries to VSTAT$(n^2/(rk^2))$, and a 1-bit sampling algorithm needs $\Omega(n^2/k^2)$
calls to 1-STAT. Search problems are bounded through the statistical dimension.
\source{Feldman, Grigorescu, Reyzin, Vempala and Xiao, J.~ACM 64(2), 2017, library key
\texttt{FGRVX17}, Theorems 2.9, 2.10, 3.2 and 5.4, Lemma 3.10, Corollary 3.12.
Accepted by astra in round 00.}
\fullproof{fact-fgrvx17-planted-clique-sq-bound}{lab/theory/spaces/computational-limits/round-00/builder.md}
\end{borrowed}

\begin{borrowed}[When statistical queries bind gradient training]\blabel{fact-sq-simulates-minibatch-sgd}\status{borrowed}
A mini-batch SGD method with batch size times squared precision above a logarithmic
factor is simulated by an SQ method, so SQ lower bounds bind it, under worst-case
rounding. At fine precision, constructed models simulate PAC learning, so a generic
SQ bound does not cover the whole fine-precision class.
\source{Abbe, Kamath, Malach, Sandon and Srebro, NeurIPS 2021, library key
\texttt{AKMSS21}, Theorems 1a and 1c. Accepted by astra in round 01.}
\fullproof{fact-sq-simulates-minibatch-sgd}{lab/theory/spaces/computational-limits/round-01/builder.md}
\end{borrowed}

\begin{borrowed}[Random expansive ReLU layers invert by linear programs]\blabel{fact-lei-random-expansive-layers-invertible}\status{borrowed}
If every layer of a ReLU network has i.i.d.\ Gaussian weights and expands its input
by a factor at least $2.1$, then for any input the input is recovered from the
output by layer-wise linear equations with probability $1-e^{-\Omega(k)}$ over the
weights. Deciding whether a preimage exists is NP-complete in general.
\source{Lei, Jalal, Dhillon and Dimakis, NeurIPS 2019, library key \texttt{LJDD19},
Assumption 1, Lemma 1, Theorems 1 to 3. Accepted by astra in round 00.}
\fullproof{fact-lei-random-expansive-layers-invertible}{lab/theory/spaces/computational-limits/round-00/builder.md}
\end{borrowed}

\begin{borrowed}[Homogeneous CLWE is a hard hidden-direction law]\blabel{fact-clwe-hardness}\status{borrowed}
For $\beta\in(0,1)$, $\gamma\ge2\sqrt d$ and $\gamma/\beta$ polynomial, distinguishing the
pancake law $H_{w,\beta,\gamma}$ from the Gaussian is as hard as a worst-case lattice
problem for quantum polynomial time. Its covariance is $I/(2\pi)$ off $w$ and
$1/(2\pi)-\Delta$ along $w$ with $0\le\Delta\le3(\gamma^2+1)e^{-\pi\gamma^2}$, and the law
determines $w$ only up to sign.
\source{Bruna, Regev, Song and Tang, STOC 2021, library key \texttt{BRST21},
Definitions 2.19 and 2.21, Corollary 4.2, Lemma 7.2, Corollary 7.3, Theorem 7.5,
Corollary 9.4. The covariance clause is proven in the round file. Accepted by argon
in round 02.}
\fullproof{fact-clwe-hardness}{lab/theory/spaces/computational-limits/round-02/builder.md}
\end{borrowed}

\begin{borrowed}[Pairwise correlations of pancake laws]\blabel{fact-hclwe-sq-correlation-bound}\status{borrowed}
Pancake laws with nearly orthogonal hidden directions have pairwise correlations,
relative to the Gaussian, exponentially small in $\gamma^2$, which gives SQ lower
bounds for detecting and finding the direction.
\source{BRST21, Proposition 8.3 and Theorem 8.1. Accepted by astra in round 01.}
\fullproof{fact-hclwe-sq-correlation-bound}{lab/theory/spaces/computational-limits/round-01/builder.md}
\end{borrowed}

\begin{borrowed}[NGCA is SQ-hard under moment matching]\blabel{fact-ngca-sq-lower-bound}\status{borrowed}
Non-Gaussian component analysis testing is hard for SQ algorithms when the hidden
law approximately matches the Gaussian moments to a given degree, with no
chi-square condition. The bound covers SQ algorithms only. Lattice reduction solves
some noiseless moment-matched instances with $O(n)$ samples.
\source{Diakonikolas, Kane, Ren and Sun, NeurIPS 2023, library key \texttt{DKRS23},
Theorem 1.5, Definitions 1.2 and 1.3. Accepted by astra in round 01.}
\fullproof{fact-ngca-sq-lower-bound}{lab/theory/spaces/computational-limits/round-01/builder.md}
\end{borrowed}

\begin{borrowed}[The price of a label-relevant direction]\blabel{fact-generative-exponent-sq-recovery-bound}\status{borrowed}
For Gaussian single-index models with generative exponent $k^*$, SQ algorithms need
of order $d^{k^*/2}$ samples to find the direction, and a sample-based algorithm with
a label denoiser matches it. The lower bound measures detection first.
\source{Damian, Pillaud-Vivien, Lee and Bruna, COLT 2024, library key \texttt{DPLB24},
Theorem 3.2, Theorem 4.1, Remark 3.3, Lemmas 2.3, E.2 and E.3. Accepted by astra in
round 01.}
\fullproof{fact-generative-exponent-sq-recovery-bound}{lab/theory/spaces/computational-limits/round-01/builder.md}
\end{borrowed}

\begin{borrowed}[Orthogonal equivariance of training]\blabel{fact-orthogonal-equivariance-of-training}\status{borrowed}
A rotation-invariant initialization with an equivariant update gives an orthogonally
equivariant learner. Gradient descent from i.i.d.\ Gaussian initialization is
orthogonally equivariant, momentum included, while Adam and AdaGrad are equivariant
only under permutations and sign flips.
\source{Li, Zhang and Arora, arXiv 2010.08515, library key \texttt{LZA21},
Definition 3.2, Theorem C.1, Remarks C.1 and C.2, Corollary C.2, Table 1. Accepted
by astra in round 01.}
\fullproof{fact-orthogonal-equivariance-of-training}{lab/theory/spaces/computational-limits/round-01/builder.md}
\end{borrowed}

\begin{borrowed}[Narcissus resiliency and one-wayness]\blabel{fact-cohen-narcissus-equals-one-wayness}\status{borrowed}
Reconstruction security measured against the baseline of a fresh sample from the
same law, Narcissus resiliency, is equivalent, for the preimage relation of an
efficiently computable well-spread function, to that function being one-way.
\source{Cohen, Kaplan, Mansour, Moran, Nissim, Stemmer and Tsfadia, ITCS 2025,
library key \texttt{CKMMNST25}, Definitions 1.6, C.1 and C.4, Theorem C.5. Accepted
by astra in round 00.}
\fullproof{fact-cohen-narcissus-equals-one-wayness}{lab/theory/spaces/computational-limits/round-00/builder.md}
\end{borrowed}

\begin{borrowed}[The lattice isomorphism problem has an average case]\blabel{fact-lip-average-case}\status{borrowed}
The lattice isomorphism problem in quadratic-form language asks, given $Q'$ in the
class of $Q$, for $U\in GL_n(\mathbb Z)$ with $Q'=U^\top QU$. For
$s\ge\lVert B_Q^*\rVert\sqrt{\ln(2n+4)/\pi}$ exact discrete Gaussian samples of
$D_{Q,s,c}$ are drawn in expected polynomial time. The Gaussian form distribution
$D_s([Q])$ draws $\lceil2n/C\rceil$ vectors from $D_{Q,s}$ until they span and returns
the form of the extracted basis, which depends only on the class. The average-case
problem asks for an isomorphism from $Q$ to $Q'\sim D_s([Q])$. An average-case oracle at
width $s\ge2^{\Theta(n)}\lambda_n(Q)$ solves the worst case with polynomial overhead, and
at $s\ge\lambda_n(Q)$ with an SIVP oracle.
\source{Ducas and van Woerden, EUROCRYPT 2022, library key \texttt{DvW22}, Definitions
2.1, 2.2, 3.3 and 3.7, Lemmas 2.9, 3.4, 3.6, 3.9 and 3.10, read in the CWI copy, PDF
not placed. Accepted by sol in round 03.}
\fullproof{fact-lip-average-case}{lab/theory/spaces/computational-limits/round-03/builder.md}
\end{borrowed}

\begin{borrowed}[The price of a rotation of $\mathbb Z^n$]\blabel{fact-zlip-best-algorithms}\status{borrowed}
For ZLIP, the lattice isomorphism problem with $Q=I_n$, a specialization of slide
reduction returns an orthogonal basis of a lattice isometric to $\mathbb Z^n$ after
$O(n^2\log\max_i\lVert b_i\rVert)$ iterations, each with an SVP oracle in dimension at
most $n/2+1$. The best provable algorithm costs $2^{n/2+o(n)}$ and the best heuristic
one $2^{0.146n+o(n)}$, as Duc24 reports them, and no polynomial-time algorithm is known.
\usesstatement{fact-lip-average-case}
\source{Ducas, Designs, Codes and Cryptography 2024, library key \texttt{Duc24}, Theorems
1 and 2, Sections 1 and 1.1, citing Bennett, Ganju, Peetathawatchai and
Stephens-Davidowitz (EUROCRYPT 2023) for the provable cost and Ducas, Postlethwaite,
Pulles and van Woerden (ASIACRYPT 2022) for the heuristic one, neither read. Accepted by
sol in round 03.}
\fullproof{fact-zlip-best-algorithms}{lab/theory/spaces/computational-limits/round-03/builder.md}
\end{borrowed}

\begin{borrowed}[LLL reads integer labels off a Gaussian background]\blabel{fact-lll-recovers-discrete-hidden-labels}\status{borrowed}
From $d+1$ samples $z_i\sim\mathcal N(ax_iu,\Sigma)$ with integer labels $x_i$, a known
spacing $a$, a unit $u$ and a $u$-weakly separable $\Sigma$, LLL on the coefficients of
the linear dependency of the samples returns in polynomial time, with probability
$1-e^{-\Omega(d)}$, an integer vector rationally proportional to the labels, provided
some label among $x_2,\dots,x_{d+1}$ is nonzero. Noiseless homogeneous CLWE with labels
from the discrete Gaussian on $(1/\gamma)\mathbb Z$ is solved this way for $\gamma$
polynomial.
\source{Zadik, Song, Wein and Bruna, COLT 2022, library key \texttt{ZSWB22}, Assumption
3.1, Model 3.2, Theorem 3.3, Lemma 3.4, Model 4.3 and Corollary 4.4. Theorem 3.3 fails
when $x_2,\dots,x_{d+1}$ all vanish, which also affects its Lemma 6.3, and the output is
rationally proportional to the labels and not an integer multiple (sol's corrections,
round 03). The paper states that inverse-polynomial noise breaks the method. Accepted
by sol in round 04.}
\fullproof{fact-lll-recovers-discrete-hidden-labels}{lab/theory/spaces/computational-limits/round-04/builder.md}
\end{borrowed}

\begin{borrowed}[Polynomially many discrete Gaussian samples generate the lattice]\blabel{fact-discrete-gaussian-samples-generate}\status{borrowed}
For a lattice of rank $n$ with determinant at least one and a width $s$ at least the
length of a shortest basis, $n^2+n\log(s\sqrt n)(n+20\log\log(s\sqrt n))$ independent
discrete Gaussian samples of width $s$ fail to generate the lattice with probability
$2^{-\Omega(n)}$. For $\mathbb Z^n$ the shortest basis has length one.
\source{Haviv and Regev, SODA 2014 as sol gives the venue, arXiv 1311.0366, library key
\texttt{HR14}, Lemma 5.4 with Lemma 5.1, PDF not placed. Accepted by sol in round 04.}
\fullproof{fact-discrete-gaussian-samples-generate}{lab/theory/spaces/computational-limits/round-04/builder.md}
\end{borrowed}

\section{From the space statistical-limits}

\begin{borrowed}[With the covariance known the sample price falls]\blabel{fact-hsu-gap-prices-covariance}\status{borrowed}
For identity-covariance Gaussians, $n\ge4$, $m\ge2$ and $d<c(n+n^2\rho^2+n^2/m)$, no
sample-based membership attack against the noisy mean exists, so at the critical
dimension an attack needs $m\gtrsim\min\{n,1/\rho^2\}$ reference samples. The price of
\cref{fact-hsu-membership-sample-gap} at that dimension is the price of not knowing the
covariance.
\source{HSU25, Theorem 4.1 and Remark 4.2, with $\rho>0$. Accepted by astra and argon.}
\fullproof{fact-hsu-gap-prices-covariance}{lab/theory/spaces/statistical-limits/round-01/builder.md}
\end{borrowed}

\begin{borrowed}[Differential privacy bounds informed membership advantage]\blabel{fact-yeom-dp-bounds-membership-advantage}\status{borrowed}
An $\varepsilon$-DP learner gives membership advantage at most $e^\varepsilon-1$ against
every adversary of the membership experiment, which is handed the data distribution.
\source{Yeom, Giacomelli, Fredrikson and Jha, IEEE CSF 2018, library key \texttt{YGFJ18},
Experiment 1, Definition 4, Theorem 1. Accepted by astra.}
\fullproof{fact-yeom-dp-bounds-membership-advantage}{lab/theory/spaces/statistical-limits/round-01/builder.md}
\end{borrowed}

\begin{borrowed}[A loss threshold gains the generalization error]\blabel{fact-yeom-generalization-error-gives-advantage}\status{borrowed}
For a loss bounded by $B$, the adversary that outputs one with probability
$\ell(A_S,z)/B$ has advantage equal to the average generalization error divided by
$B$, with one model evaluation and no sample or training run.
\source{YGFJ18, Theorem 2 and Definition 3. Accepted by astra and argon.}
\fullproof{fact-yeom-generalization-error-gives-advantage}{lab/theory/spaces/statistical-limits/round-01/builder.md}
\end{borrowed}

\begin{borrowed}[The LiRA game excludes members from the OUT draw]\blabel{fact-lira-low-fpr-yeom-game}\status{borrowed}
The membership game of LiRA draws its OUT point from the distribution conditioned to
lie outside the training set, with the bit renamed against Yeom's experiment, so the two
games agree when that event has probability one.
\source{CCNSTT22, Definition 1 in \S III-A. Accepted by astra and argon.}
\fullproof{fact-lira-low-fpr-yeom-game}{lab/theory/spaces/statistical-limits/round-01/builder.md}
\end{borrowed}

\begin{borrowed}[Reconstruction robustness and its baseline]\blabel{fact-balle-reconstruction-robustness-definition}\status{borrowed}
A mechanism is $(\eta,\gamma)$-reconstruction-robust for a prior and an error when every
deterministic reconstructor from the released model, against an informed adversary
holding all other records, succeeds within $\eta$ with probability at most $\gamma$.
The baseline $\kappa(\eta)$ is the best chance of a fixed guess under the prior.
\source{Balle, Cherubin and Hayes, IEEE S\&P 2022, library key \texttt{BCH22},
Definitions 1 and 2, \S VII-B. Accepted by astra and argon.}
\fullproof{fact-balle-reconstruction-robustness-definition}{lab/theory/spaces/statistical-limits/round-01/builder.md}
\end{borrowed}

\begin{borrowed}[Differential privacy bounds reconstruction by the baseline]\blabel{fact-balle-dp-bounds-reconstruction}\status{borrowed}
An $(\alpha,\varepsilon)$-RDP mechanism is reconstruction-robust with
$\gamma=(\kappa e^\varepsilon)^{(\alpha-1)/\alpha}$, an $\varepsilon$-DP one with
$\gamma=\kappa e^\varepsilon$, and a $\rho$-zCDP one with
$\gamma=\exp(-(\sqrt{\log(1/\kappa)}-\sqrt\rho)^2)$. Conversely, robustness against
every two-point prior gives differential privacy.
\source{BCH22, Theorem 2, Corollaries 3 and 4, Propositions 6 and 7, Theorem 5.
Accepted by astra and argon.}
\fullproof{fact-balle-dp-bounds-reconstruction}{lab/theory/spaces/statistical-limits/round-01/builder.md}
\end{borrowed}

\begin{borrowed}[How many outputs a reconstruction attacker must see]\blabel{fact-reconstruction-query-lower-bounds}\status{borrowed}
For every $(\varepsilon,\delta)$-DP learner on a compact metric space, an adversary that
sees $n$ independent outputs has worst-case squared reconstruction error of order at
least the squared diameter times $(2(1-\delta)/(e^\varepsilon+1))^n$. The claimed matching
upper bound and its R\'enyi extension are not established.
\source{Mukherjee and Lokam, arXiv 2303.16372v6, library key \texttt{ML24}, Theorem 9. Theorem 10 rests on an adversary that uses the true target, and
the proof of Lemma 16 bounds the total variation from the wrong side. Accepted by argon
in round 02.}
\fullproof{fact-reconstruction-query-lower-bounds}{lab/theory/spaces/statistical-limits/round-01/builder.md,lab/theory/spaces/statistical-limits/round-02/builder.md}
\end{borrowed}

\begin{borrowed}[Reconstruction thresholds for subset sums]\blabel{fact-dinur-nissim-reconstruction-thresholds}\status{borrowed}
One reconstruction theorem uses all $2^n$ subset-sum answers and exponential time,
and another uses $n(\log n)^2$ random answers and a linear program. The two results change the number of answers as
well as the running time.
\source{Dinur and Nissim, PODS 2003, library key \texttt{DN03}, Definitions 1 to 4,
Theorems 2, 3 and 5. Accepted by astra and argon.}
\fullproof{fact-dinur-nissim-reconstruction-thresholds}{lab/theory/spaces/statistical-limits/round-01/builder.md}
\end{borrowed}

\begin{borrowed}[Mutual information bounds the gain over the prior guess]\blabel{fact-pac-privacy-mutual-information-bound}\status{borrowed}
For a data law known to the adversary, an adversary whose failure probability is at
most the optimal prior failure $\delta_o$ satisfies
$\mathrm{kl}(\delta\|\delta_o)\le I(X;M(X))$, and every adversary has
$\delta_o-\delta\le\sqrt{I(X;M(X))/2}$. Side knowledge conditions both sides.
\source{Xiao and Devadas, CRYPTO 2023, library key \texttt{XD23}, Definition 3,
Theorem 1, Corollary 1, Remark 1. Accepted by astra and argon.}
\fullproof{fact-pac-privacy-mutual-information-bound}{lab/theory/spaces/statistical-limits/round-01/builder.md}
\end{borrowed}

\begin{borrowed}[A Gaussian release has information set by its covariance]\blabel{fact-pac-privacy-gaussian-release-bound}\status{borrowed}
For a deterministic $M$ with finite second moment and independent Gaussian noise of
positive definite covariance $\Sigma_B$,
$I(X;M(X)+B)\le\frac12\log\det(I+\Sigma_{M(X)}\Sigma_B^{-1})$, and an algorithm that
evaluates $M$ on independent draws calibrates the noise with confidence.
\source{XD23, Theorems 3 and 4, the deterministic specialization of \S 4.2. Accepted by
astra and argon.}
\fullproof{fact-pac-privacy-gaussian-release-bound}{lab/theory/spaces/statistical-limits/round-01/builder.md}
\end{borrowed}

\begin{borrowed}[An accurate representation leaks a correlated attribute]\blabel{fact-zhao-accuracy-forces-attribute-leakage}\status{borrowed}
When the target and a sensitive attribute are correlated, every representation on
which some classifier is accurate lets a measurable classifier infer the attribute
with an advantage that the accuracy forces, for attribute groups of positive
probability.
\source{Zhao, Chi, Tian and Gordon, NeurIPS 2020, library key \texttt{ZCTG20}, Theorems 3.2
and 3.3.
Accepted by astra and argon.}
\fullproof{fact-zhao-accuracy-forces-attribute-leakage}{lab/theory/spaces/statistical-limits/round-01/builder.md}
\end{borrowed}

\begin{borrowed}[A learned reconstructor costs thousands of shadow runs]\blabel{fact-balle-reconnn-priced-in-shadow-runs}\status{borrowed}
A reconstructor network trained on pairs of shadow models and their target records
beats the nearest-neighbour oracle from about 8{,}000 shadow models on MNIST and
10{,}000 on CIFAR-10 when the initialization is known, and fails on MNIST with a fresh
initialization per model.
\source{BCH22, \S IV-B, \S VI-B, Fig.~5, Table III. Empirical. One attack in two
experiments, not a statement that reconstruction without the initialization is
impossible. Accepted by astra and argon.}
\fullproof{fact-balle-reconnn-priced-in-shadow-runs}{lab/theory/spaces/statistical-limits/round-01/builder.md}
\end{borrowed}

\begin{borrowed}[The sample covariance of sub-gaussian rows concentrates]\blabel{fact-subgaussian-sample-covariance-concentration}\status{borrowed}
Let $A$ be an $N\times d$ matrix with independent, mean-zero, isotropic, sub-gaussian
rows of $\psi_2$ norm at most $K$. For every $u\ge0$, with probability at least
$1-2e^{-u^2}$, $\lVert A^\top A/N-I_d\rVert\le K^2\max(\delta,\delta^2)$ with
$\delta=C(\sqrt{d/N}+u/\sqrt N)$ for an absolute constant $C$.
\source{BRST21 Lemma 7.4, which quotes Vershynin, High-Dimensional Probability (2018),
Theorem 4.6.1, a book not in the library. fable's critique of the debate space notes,
from memory, that the book places $K^2$ inside $\delta$, so the quadratic branch there
carries $K^4$. Every use in this booklet stays in the linear branch, where the two forms
agree up to the absolute constant. Accepted by sol in rounds 03 to 05.}
\fullproof{fact-subgaussian-sample-covariance-concentration}{lab/theory/spaces/statistical-limits/round-03/builder.md}
\end{borrowed}

\section{From the space hard-problem-mappings}

\begin{borrowed}[Label post-processing never lowers the generative exponent]\blabel{fact-generative-exponent-post-processing-monotone}\status{borrowed}
For the class $\mathcal G$ of laws of $(Z,Y)$ with Gaussian $Z$, finite second moment
and positive $\chi^2$ information, the generative exponent $k^*$ is the least $k$ with
$\lambda_k=\lVert\mathbb E[h_k(Z)\mid Y]\rVert\ne0$, and the $\chi^2$ information is
$\sum_k\lambda_k^2$. It is the infimum over label transforms of the information
exponent of the transformed law. If $Y\to Y'$ is a Markov chain with
$\mathbb E[Y'^2]<\infty$ then $k^*(Z,Y')\ge k^*(Z,Y)$, invertible label maps leave it
unchanged, and Gaussian label noise on a deterministic link preserves it.
\notation{$k^*$}{generative exponent of a single-index law, the least Hermite degree with a nonzero conditional moment}
\source{Damian, Pillaud-Vivien, Lee and Bruna, COLT 2024, library key \texttt{DPLB24},
Definitions 1.1 and 2.4, Lemma 2.3, Proposition 2.6 and the paragraph after it,
Example 2.7 (ii), Corollary 2.8, Theorem 5.2. The proof of its Theorem 3.2 takes the
diagonal correlation bound to be one where it equals the $\chi^2$ information, which
can exceed one and is infinite for deterministic labels, so a bound for a deterministic
link needs label noise (the space, from Appendix E.2). Accepted by sol in round 00.}
\fullproof{fact-generative-exponent-post-processing-monotone}{lab/theory/spaces/hard-problem-mappings/round-00/builder.md}
\end{borrowed}

\begin{borrowed}[Backdoors planted in the weight law of random features]\blabel{fact-gkvz22-backdoor-scope}\status{borrowed}
Under one-way functions every training procedure admits black-box undetectable
backdoors, built from digital signatures. White-box undetectability is shown for random
Fourier features whose weights come from sparse Gaussian pancakes under CLWE, with a
normalized head of inverse-polynomial margin and width
$m=\Theta(d\log(d/\varepsilon\delta)/\varepsilon^2)$. The second white-box construction,
random ReLU features with a sparse-PCA spike in the weights, fails at its printed
parameters (\cref{lem-sparse-pca-backdoor-width-price}).
\source{Goldwasser, Kim, Vaikuntanathan and Zamir, arXiv 2204.06974, library key
\texttt{GKVZ22}, Definition 4.6, Theorems 2.1, 2.2, 5.4 and 6.4, Lemmas 6.2 and 6.6,
Appendix A with Claim A.5 and Lemma A.3. Accepted by sol in round 01.}
\fullproof{fact-gkvz22-backdoor-scope}{lab/theory/spaces/hard-problem-mappings/round-01/builder.md}
\end{borrowed}
